% Generated by publication/build_sections.py; do not edit by hand.
\paragraph{Informal verdict.}
\textbf{C --- Substantial progress}, with \textbf{high confidence in the stated exact special cases}.

\paragraph{Lean coverage (partial).}
The published/imported generation-1 source remains the verified 3,644-line partial corpus. The latest archive reports compiling 5,366 Lean lines across Q706, Q706Ext, and Q706Second, adding general compactness and attainment, singleton boxes, global bounds and reductions, interlacing, new signed second/penultimate cells, all endpoints for n<=3, homogeneity, translation, and etaConst\_penultimate. General intermediate endpoints remain unevaluated; q706\_aggregate uses sInf/sSup and conditional known cells, and the audited n<=5 plus eta\_(6,2)=4 browser results are not in Lean.

\subsection{Audited informal answer}
\subsection{Partial result}

\subsubsection{1. Faithful restatement}

For \(a\le b\), let \(S_n([a,b])\) be the set of real symmetric \(n\times n\) matrices all of whose entries belong to ([a,b]). For a real symmetric matrix \(A\), write

\[
\lambda_1(A)\ge \lambda_2(A)\ge\cdots\ge \lambda_n(A)
\]

for its eigenvalues in decreasing order. The question asks for the lower and upper endpoints of

\[
\bigl\{\lambda_k(A):A\in S_n([a,b])\bigr\}
\]

as functions of (n,a,b). The statement implicitly assumes \(1\le k\le n\); there is no other ambiguity or apparent typo. 

Define

\[
L_k(n;a,b):=\min_{A\in S_n([a,b])}\lambda_k(A),\qquad
U_k(n;a,b):=\max_{A\in S_n([a,b])}\lambda_k(A).
\]

The set \(S_n([a,b])\) is compact and connected, and \(A\mapsto\lambda_k(A)\) is continuous. Hence the set asked for in Q706 is exactly the compact interval

\[
\boxed{\ [L_k(n;a,b),U_k(n;a,b)]\ }.
\]

I obtain complete closed formulas for \(k=1,n\), complete formulas for \(k=2\) when \(0\le a\le b\), and---by sign reversal---for \(k=n-1\) when \(a\le b\le0\). For every \(k\), I also give an exact invariant minimax characterization and explicit universal bounds.

\medskip\hrule\medskip

\subsection{2. Main results}

Throughout, put

\[
c:=\frac{a+b}{2},\qquad r:=\frac{b-a}{2},\qquad
\rho:=\max{|a|,|b|}.
\]

\subsubsection{2.1 Exact answer for the largest and smallest ordered eigenvalues}

For every \(n\ge2\),

\[
\boxed{
L_1(n;a,b)=
\begin{cases}
na, & a\ge0,\\[2mm]
a, & a\le0\le b,\\[2mm]
a-b, & b\le0,
\end{cases}}
\tag{2.1}
\]

and

\[
\boxed{
U_n(n;a,b)=
\begin{cases}
b-a, & a\ge0,\\[2mm]
b, & a\le0\le b,\\[2mm]
nb, & b\le0.
\end{cases}}
\tag{2.2}
\]

The other two outer endpoints are

\[
\boxed{
U_1(n;a,b)=
\begin{cases}
nb, & a+b\ge0,\\[2mm]
\dfrac{n(b-a)}2,
& a+b\le0,\ n\ \text{even},\\[3mm]
\dfrac{nb+\sqrt{(n^2-1)a^2+b^2}}2,
& a+b\le0,\ n\ \text{odd},
\end{cases}}
\tag{2.3}
\]

and

\[
\boxed{
L_n(n;a,b)=
\begin{cases}
na, & a+b\le0,\\[2mm]
\dfrac{n(a-b)}2,
& a+b\ge0,\ n\ \text{even},\\[3mm]
\dfrac{na-\sqrt{a^2+(n^2-1)b^2}}2,
& a+b\ge0,\ n\ \text{odd}.
\end{cases}}
\tag{2.4}
\]

The formulas agree on the boundary \(a+b=0\).

Thus the ranges of \(\lambda_1\) and \(\lambda_n\) are completely determined:

\[
{\lambda_1(A):A\in S_n([a,b])}
=[L_1(n;a,b),U_1(n;a,b)],
\]

\[
{\lambda_n(A):A\in S_n([a,b])}
=[L_n(n;a,b),U_n(n;a,b)].
\]

\medskip\hrule\medskip

\subsubsection{2.2 Exact answer for \(\lambda_2\) when the interval is nonnegative}

Assume

\[
0\le a\le b.
\]

Then

\[
\boxed{L_2(n;a,b)=a-b}
\tag{2.5}
\]

and

\[
\boxed{
U_2(n;a,b)=
\begin{cases}
\dfrac{n(b-a)}2,
& n\ \text{even},\\[3mm]
\dfrac{nb-\sqrt{(n^2-1)a^2+b^2}}2,
& n\ \text{odd}.
\end{cases}}
\tag{2.6}
\]

Consequently,

\[
\boxed{
{\lambda_2(A):A\in S_n([a,b])}
=
%
\left[
a-b,
\begin{cases}
\dfrac{n(b-a)}2,&n\ \mathrm{even},\\[2mm]
\dfrac{nb-\sqrt{(n^2-1)a^2+b^2}}2,&n\ \mathrm{odd}
\end{cases}
\right].}
\]

In particular, for \(a=0\),

\[
\boxed{U_2(n;0,b)=b\left\lfloor\frac n2\right\rfloor.}
\tag{2.7}
\]

\medskip\hrule\medskip

\subsubsection{2.3 Exact answer for \(\lambda_{n-1}\) when the interval is nonpositive}

By applying (2.5)--(2.6) to (-A), if

\[
a\le b\le0,
\]

then

\[
\boxed{U_{n-1}(n;a,b)=b-a}
\tag{2.8}
\]

and

\[
\boxed{
L_{n-1}(n;a,b)=
\begin{cases}
\dfrac{n(a-b)}2,
& n\ \text{even},\\[3mm]
\dfrac{na+\sqrt{a^2+(n^2-1)b^2}}2,
& n\ \text{odd}.
\end{cases}}
\tag{2.9}
\]

\medskip\hrule\medskip

\subsubsection{2.4 Exact invariant characterization for every \(k\)}

Let \(\mathbf 1=(1,\dots,1)^T\). For a symmetric matrix \(X=(x_{ij})\), write

\[
|X|_{1,e}:=\sum_{i,j=1}^n |x_{ij}|
\]

for its entrywise \(1\)-norm.

For a subspace \(V\subseteq\mathbb R^n\), define

\[
\mathcal D(V):=
\left\{
X\in S_n(\mathbb R):
X\succeq0,\\
\operatorname{tr}X=1,\\
\operatorname{Ran}X\subseteq V
\right\}.
\]

Then for every \(1\le k\le n\),

\[
\boxed{
U_k(n;a,b)
=
%
\max_{\substack{V\le\mathbb R^n\ \dim V=k}}
\ \min_{X\in\mathcal D(V)}
\left(
c,\mathbf1^TX\mathbf1+r|X|_{1,e}
\right).}
\tag{2.10}
\]

Putting \(q=n-k+1\),

\[
\boxed{
L_k(n;a,b)
=
%
\min_{\substack{W\le\mathbb R^n\ \dim W=q}}
\ \max_{X\in\mathcal D(W)}
\left(
c,\mathbf1^TX\mathbf1-r|X|_{1,e}
\right).}
\tag{2.11}
\]

These are exact formulas, not relaxations.

For the centered interval ([-r,r]), they simplify to

\[
\boxed{
U_k(n;-r,r)
=
%
r,\eta_{n,k},
\qquad
\eta_{n,k}:=
\max_{\dim V=k}\min_{X\in\mathcal D(V)}|X|_{1,e}.}
\tag{2.12}
\]

Furthermore,

\[
L_k(n;-r,r)=-U_{n-k+1}(n;-r,r).
\tag{2.13}
\]

The constants satisfy

\[
\eta_{n,1}=n,\qquad \eta_{n,n}=1.
\]

Evaluating \(\eta_{n,k}\) explicitly for all intermediate \(k\) is the central remaining obstruction.

\medskip\hrule\medskip

\subsubsection{2.5 General explicit bounds for all \(k\)}

For every \(1\le k\le n\),

\[
\boxed{
U_k(n;a,b)
\le
b+\rho
\sqrt{\frac{(n-k)(n-1)}{k}}}
\tag{2.14}
\]

and

\[
\boxed{
L_k(n;a,b)
\ge
a-\rho
\sqrt{\frac{(k-1)(n-1)}{n-k+1}}.}
\tag{2.15}
\]

A second family of bounds follows from the geometry of orthogonal projections. Let

\[
c_+:=\max(c,0),\qquad c_-:=\min(c,0).
\]

An elementary estimate gives

\[
\boxed{
U_k(n;a,b)
\le
\frac nk\bigl(c_++r\sqrt{k}\bigr)}
\tag{2.16}
\]

and, with \(q=n-k+1\),

\[
\boxed{
L_k(n;a,b)
\ge
\frac nq\bigl(c_- -r\sqrt q\bigr).}
\tag{2.17}
\]

A sharper known projection estimate replaces \(\sqrt s\) by

\[
\beta_s:=
\frac{s+\sqrt{s+2}}{1+\sqrt{s+2}},
\]

giving

\[
\boxed{
U_k(n;a,b)
\le
\frac nk\bigl(c_++r\beta_k\bigr)}
\tag{2.18}
\]

and

\[
\boxed{
L_k(n;a,b)
\ge
\frac n{n-k+1}
\bigl(c_- -r\beta_{,n-k+1}\bigr).}
\tag{2.19}
\]

The external input here is the theorem that every rank-\(s\) real orthogonal projection \(P\) of order \(n\) satisfies

\[
\sum_{i,j}|p_{ij}|\le \beta_s n.
\]

Its hypotheses match exactly the spectral projections used below. ([arXiv][1])

\medskip\hrule\medskip

\subsection{3. Preliminary reductions}

\subsubsection{3.1 Sign reversal}

If \(A\in S_n([a,b])\), then

\[
-A\in S_n([-b,-a])
\]

and

\[
\lambda_k(-A)=-\lambda_{n-k+1}(A).
\]

Therefore

\[
\boxed{
L_k(n;a,b)
=
%
-U_{n-k+1}(n;-b,-a).}
\tag{3.1}
\]

This will repeatedly halve the work.

\medskip\hrule\medskip

\subsubsection{3.2 The diagonal may be fixed at an endpoint}

Suppose (A') is obtained from \(A\) by increasing some diagonal entries. Then

\[
A'-A\succeq0.
\]

By the Courant--Fischer principle,

\[
\lambda_k(A')\ge\lambda_k(A)
\]

for every \(k\). Consequently:

\begin{itemize}
\item in computing \(U_k(n;a,b)\), all diagonal entries may be taken equal to \(b\);
\item in computing \(L_k(n;a,b)\), all diagonal entries may be taken equal to \(a\).
\end{itemize}

This observation is used throughout.

\medskip\hrule\medskip

\subsection{4. Proof of the exact outer formulas}

\subsubsection{4.1 Maximizing \(\lambda_1\)}

For a fixed unit vector \(x=(x_i)\),

\[
x^TAx=\sum_{i,j}a_{ij}x_ix_j.
\]

Each coefficient \(a_{ij}\) can be chosen independently, subject only to symmetry. It is optimal to choose \(b\) when \(x_ix_j\ge0\) and \(a\) when \(x_ix_j<0\). Thus

\[
\max_{A\in S_n([a,b])}x^TAx
=
%
c\left(\sum_i x_i\right)^2
+r\left(\sum_i|x_i|\right)^2.
\tag{4.1}
\]

Since the two maximizations commute,

\[
U_1(n;a,b)
=
%
\max_{|x|_2=1}
\left[
c(\mathbf1^Tx)^2+r|x|_1^2
\right].
\tag{4.2}
\]

\paragraph{Case \(a+b\ge0\)}

Here \(c\ge0\). Since

\[
|\mathbf1^Tx|\le|x|_1\le\sqrt n,
\]

we obtain

\[
c(\mathbf1^Tx)^2+r|x|_1^2
\le(c+r)|x|_1^2
=b|x|_1^2
\le nb.
\]

Equality is attained for

\[
x=\frac1{\sqrt n}\mathbf1,\qquad A=bJ_n.
\]

Hence

\[
U_1=nb.
\]

\paragraph{Case \(a+b<0\)}

In this case \(a<0\) and \(a^2>b^2\).

For a sign vector \(\varepsilon\in{\pm1}^n\), define the endpoint matrix

\[
B_\varepsilon=(b_{ij}),\qquad
b_{ij}=
\begin{cases}
b,&\varepsilon_i=\varepsilon_j,\\
a,&\varepsilon_i\ne\varepsilon_j.
\end{cases}
\]

For each \(x\), choosing \(\varepsilon_i=\operatorname{sgn}x_i\) gives the right-hand side of (4.1), so

\[
U_1=\max_{\varepsilon}\lambda_1(B_\varepsilon).
\tag{4.3}
\]

Suppose \(\varepsilon\) has \(p\) plus signs and \(q=n-p\) minus signs. Up to permutation,

\[
B_\varepsilon=
\begin{pmatrix}
bJ_p&aJ_{p,q}\\
aJ_{q,p}&bJ_q
\end{pmatrix}.
\tag{4.4}
\]

It has \(n-2\) zero eigenvalues, and its remaining two eigenvalues are those of

\[
\begin{pmatrix}
bp&a\sqrt{pq}\\
a\sqrt{pq}&bq
\end{pmatrix}.
\]

Thus

\[
\lambda_1(B_\varepsilon)
=
%
\frac{
bn+
\sqrt{b^2(p-q)^2+4a^2pq}
}{2}.
\tag{4.5}
\]

Since

\[
b^2(p-q)^2+4a^2pq
=
%
a^2n^2+(b^2-a^2)(p-q)^2
\]

and \(b^2-a^2<0\), the expression is maximized when (|p-q|) is as small as possible.

If \(n\) is even, \(p=q=n/2\), giving

\[
U_1=\frac{n(b-a)}2.
\]

If \(n\) is odd, \(|p-q|=1\) and \(4pq=n^2-1\), giving

\[
U_1=
\frac{nb+\sqrt{(n^2-1)a^2+b^2}}2.
\]

This proves (2.3).

\medskip\hrule\medskip

\subsubsection{4.2 Minimizing \(\lambda_1\)}

The map \(A\mapsto\lambda_1(A)\) is convex and invariant under simultaneous permutation of rows and columns.

Given \(A\), average \(PAP^T\) over all permutation matrices \(P\). The result has the form

\[
\overline A=(d-e)I+eJ,
\]

where \(d,e\in[a,b]\). Convexity gives

\[
\lambda_1(\overline A)\le\lambda_1(A).
\]

Hence the minimum is attained among such two-parameter matrices.

Their eigenvalues are

\[
d+(n-1)e
\]

on \(\mathbb R\mathbf1\), and

\[
d-e
\]

with multiplicity \(n-1\). Increasing \(d\) increases every eigenvalue, so \(d=a\). We must therefore minimize

\[
f(e)=\max{a+(n-1)e,\ a-e},
\qquad e\in[a,b].
\tag{4.6}
\]

If \(e\ge0\), the first term is larger. If \(e\le0\), the second is larger.

\begin{itemize}
\item If \(a\ge0\), the minimum occurs at \(e=a\), and equals \(na\).
\item If \(a\le0\le b\), the minimum occurs at \(e=0\), and equals \(a\).
\item If \(b\le0\), the minimum occurs at \(e=b\), and equals \(a-b\).
\end{itemize}

This proves (2.1).

The extremizing matrices may be chosen respectively as

\[
aJ,\qquad aI,\qquad (a-b)I+bJ.
\]

\medskip\hrule\medskip

\subsubsection{4.3 The formulas for \(L_n\) and \(U_n\)}

Apply the sign-reversal identity (3.1) to the formulas just proved:

\[
L_n(n;a,b)=-U_1(n;-b,-a),
\]

\[
U_n(n;a,b)=-L_1(n;-b,-a).
\]

This gives exactly (2.2) and (2.4).

For example, when \(a+b>0\), a minimizing matrix for \(\lambda_n\) is obtained from a balanced partition \(p+q=n\):

\[
\begin{pmatrix}
aJ_p&bJ_{p,q}\\
bJ_{q,p}&aJ_q
\end{pmatrix}.
\]

\medskip\hrule\medskip

\subsection{5. Exact \(\lambda_2\)-range for nonnegative intervals}

Assume \(0\le a\le b\).

\subsubsection{5.1 The lower endpoint}

In computing \(L_2\), set every diagonal entry equal to \(a\).

For any two indices \(i\ne j\), the corresponding principal \(2\times2\) submatrix is

\[
\begin{pmatrix}
a&t\\
t&a
\end{pmatrix},
\qquad t\in[a,b].
\]

Its smaller eigenvalue is \(a-t\ge a-b\).

By Courant--Fischer, restricting the two-dimensional maximizing subspace to the coordinate plane generated by \(e_i,e_j\) gives the interlacing inequality

\[
\lambda_2(A)\ge a-t\ge a-b.
\]

Conversely,

\[
A=(a-b)I+bJ
\]

belongs to \(S_n([a,b])\) and has eigenvalues

\[
a+(n-1)b,\qquad a-b,\dots,a-b.
\]

Hence

\[
L_2=a-b.
\]

\medskip\hrule\medskip

\subsubsection{5.2 The upper endpoint when \(a>0\)}

By diagonal monotonicity, take \(a_{ii}=b\).

Because every entry of \(A\) is strictly positive, the top eigenvalue is simple and admits an eigenvector with all coordinates strictly positive. Indeed, replacing a vector by its coordinatewise absolute value does not decrease its Rayleigh quotient, and strict positivity of the entries gives strict improvement for a vector having both signs.

Let \(x\) be an eigenvector for \(\lambda=\lambda_2(A)\), orthogonal to the positive top eigenvector. Hence \(x\) has both positive and negative coordinates.

Set

\[
P={i:x_i>0},\qquad N={i:x_i<0},
\]

\[
p=|P|,\qquad q=|N|,
\]

and

\[
U=\sum_{i\in P}x_i>0,\qquad
V=-\sum_{j\in N}x_j>0.
\]

For \(i\in P\),

\[
\begin{aligned}
\lambda x_i
&=\sum_{j\in P}a_{ij}x_j
+\sum_{j\in N}a_{ij}x_j\\
&\le bU-aV.
\end{aligned}
\]

Summing over \(i\in P\) gives

\[
\lambda U\le p(bU-aV).
\tag{5.1}
\]

Similarly, for \(j\in N\),

\[
\lambda(-x_j)\le bV-aU,
\]

and hence

\[
\lambda V\le q(bV-aU).
\tag{5.2}
\]

If \(\lambda\le0\), the desired positive upper bound is immediate. Assume \(\lambda>0\). Equations (5.1)--(5.2) imply

\[
(bp-\lambda)U\ge apV,
\]

\[
(bq-\lambda)V\ge aqU.
\]

Multiplying and cancelling \(UV>0\),

\[
(bp-\lambda)(bq-\lambda)\ge a^2pq.
\tag{5.3}
\]

Moreover \(bp-\lambda>0\) and \(bq-\lambda>0\). Therefore \(\lambda\) is no larger than the smaller root of

\[
(bp-t)(bq-t)-a^2pq=0.
\]

Thus

\[
\lambda\le
t_-(p,q):=
\frac{
b(p+q)-
\sqrt{b^2(p-q)^2+4a^2pq}
}{2}.
\tag{5.4}
\]

Write \(m=p+q\le n\). Since

\[
b^2(p-q)^2+4a^2pq
=
%
a^2m^2+(b^2-a^2)(p-q)^2,
\]

for fixed \(m\) the right-hand side of (5.4) is maximized by making (|p-q|) as small as possible.

Define

\[
\phi_m:=
\begin{cases}
\dfrac{m(b-a)}2,&m\ \text{even},\\[3mm]
\dfrac{mb-\sqrt{(m^2-1)a^2+b^2}}2,&m\ \text{odd}.
\end{cases}
\]

Then \(t_-(p,q)\le\phi_m\).

It remains to note that \(\phi_m\) is nondecreasing in \(m\). For even \(m\),

\[
\phi_{m+1}-\phi_m
=
%
\frac{
b+ma-\sqrt{b^2+a^2m(m+2)}
}{2}\ge0,
\]

because

\[
(b+ma)^2-\bigl(b^2+a^2m(m+2)\bigr)
=2am(b-a)\ge0.
\]

For odd \(m\),

\[
\phi_{m+1}-\phi_m
=
%
\frac{
b-(m+1)a+\sqrt{a^2(m^2-1)+b^2}
}{2}\ge0.
\]

If \((m+1)a\le b\), this is immediate. Otherwise it follows by squaring, since

\[
a^2(m^2-1)+b^2-\bigl((m+1)a-b\bigr)^2
=
%
2a(m+1)(b-a)\ge0.
\]

Consequently,

\[
\lambda_2(A)\le\phi_n.
\]

\medskip\hrule\medskip

\subsubsection{5.3 The case \(a=0\)}

For \(0<\varepsilon<1\), put

\[
A_\varepsilon=(1-\varepsilon)A+\varepsilon bJ.
\]

Then \(A_\varepsilon\) has diagonal \(b\) and all its entries lie in \([\varepsilon b,b]\). The preceding result gives

\[
\lambda_2(A_\varepsilon)\le\phi_n(\varepsilon b,b).
\]

Letting \(\varepsilon\downarrow0\) and using continuity of the eigenvalues yields the formula for \(a=0\).

\medskip\hrule\medskip

\subsubsection{5.4 Attainment}

Take a partition with

\[
p=\left\lfloor\frac n2\right\rfloor,\qquad
q=\left\lceil\frac n2\right\rceil
\]

and define

\[
A_{p,q}:=
\begin{pmatrix}
bJ_p&aJ_{p,q}\\
aJ_{q,p}&bJ_q
\end{pmatrix}.
\tag{5.5}
\]

It has \(n-2\) zero eigenvalues, while its two remaining eigenvalues are

\[
\frac{
bn\pm\sqrt{b^2(p-q)^2+4a^2pq}
}{2}.
\]

Both are nonnegative, and the smaller one is therefore \(\lambda_2(A_{p,q})\). It equals exactly the right-hand side of (2.6).

This proves the complete nonnegative-interval result.

\medskip\hrule\medskip

\subsection{6. Exact minimax characterization for arbitrary \(k\)}

We now prove (2.10)--(2.11).

\subsubsection{6.1 A bilinear endpoint computation}

Let \(X\succeq0\) with \(\operatorname{tr}X=1\). In particular \(x_{ii}\ge0\).

Since

\[
\operatorname{tr}(AX)
=
%
\sum_i a_{ii}x_{ii}
+
2\sum_{i<j}a_{ij}x_{ij},
\]

entrywise maximization gives

\[
\max_{A\in S_n([a,b])}\operatorname{tr}(AX)
=
%
c,\mathbf1^TX\mathbf1+r|X|_{1,e}.
\tag{6.1}
\]

Similarly,

\[
\min_{A\in S_n([a,b])}\operatorname{tr}(AX)
=
%
c,\mathbf1^TX\mathbf1-r|X|_{1,e}.
\tag{6.2}
\]

\medskip\hrule\medskip

\subsubsection{6.2 Upper endpoint}

Courant--Fischer gives

\[
\lambda_k(A)
=
%
\max_{\dim V=k}
\min_{\substack{x\in V\|x|=1}}
x^TAx.
\]

For fixed \(V\),

\[
\min_{\substack{x\in V\|x|=1}}x^TAx
=
%
\min_{X\in\mathcal D(V)}\operatorname{tr}(AX),
\]

because a linear functional on the compact convex set \(\mathcal D(V)\) attains its minimum at a rank-one density matrix \(xx^T\).

Therefore

\[
U_k
=
%
\max_{\dim V=k}
\max_{A\in S_n([a,b])}
\min_{X\in\mathcal D(V)}
\operatorname{tr}(AX).
\]

For fixed \(V\), the sets of admissible \(A\) and \(X\) are compact and convex, and the function \((A,X)\mapsto\operatorname{tr}(AX)\) is bilinear. The finite-dimensional minimax theorem thus gives

\[
\max_A\min_X\operatorname{tr}(AX)
=
%
\min_X\max_A\operatorname{tr}(AX).
\]

Using (6.1) proves (2.10).

\medskip\hrule\medskip

\subsubsection{6.3 Lower endpoint}

Let \(q=n-k+1\). The second form of Courant--Fischer is

\[
\lambda_k(A)
=
%
\min_{\dim W=q}
\max_{\substack{x\in W\|x|=1}}
x^TAx.
\]

As above,

\[
\max_{\substack{x\in W\|x|=1}}x^TAx
=
%
\max_{X\in\mathcal D(W)}\operatorname{tr}(AX).
\]

Thus

\[
L_k
=
%
\min_{\dim W=q}
\min_A
\max_{X\in\mathcal D(W)}
\operatorname{tr}(AX).
\]

Applying minimax in the opposite order and then (6.2) gives (2.11).

\medskip\hrule\medskip

\subsection{7. Ky Fan envelope and projection bounds}

Let \(\mathcal P_s(n)\) denote the rank-\(s\) orthogonal projections in \(\mathbb R^n\).

The Ky Fan maximum principle gives

\[
\sum_{i=1}^s\lambda_i(A)
=
%
\max_{P\in\mathcal P_s(n)}\operatorname{tr}(PA).
\]

Consequently,

\[
\boxed{
\max_{A\in S_n([a,b])}
\sum_{i=1}^s\lambda_i(A)
=
%
\max_{P\in\mathcal P_s(n)}
\left(
c,\mathbf1^TP\mathbf1+r|P|_{1,e}
\right).}
\tag{7.1}
\]

Similarly,

\[
\boxed{
\min_{A\in S_n([a,b])}
\sum_{i=n-s+1}^n\lambda_i(A)
=
%
\min_{P\in\mathcal P_s(n)}
\left(
c,\mathbf1^TP\mathbf1-r|P|_{1,e}
\right).}
\tag{7.2}
\]

Since

\[
s\lambda_s(A)\le\sum_{i=1}^s\lambda_i(A),
\]

equation (7.1) yields an upper bound for \(U_s\). Likewise, with \(q=n-k+1\),

\[
\sum_{i=k}^n\lambda_i(A)\le q\lambda_k(A),
\]

so (7.2) yields a lower bound for \(L_k\).

For an orthogonal projection \(P\),

\[
0\le\mathbf1^TP\mathbf1=|P\mathbf1|^2\le n,
\tag{7.3}
\]

and, by Cauchy--Schwarz,

\[
\begin{aligned}
|P|_{1,e}
&\le n\left(\sum_{i,j}p_{ij}^2\right)^{1/2}\\
&=n\bigl(\operatorname{tr}P^2\bigr)^{1/2}\\
&=n\sqrt s.
\end{aligned}
\tag{7.4}
\]

Equations (2.16)--(2.17) follow. Replacing (7.4) by the sharper projection theorem stated after (2.19) proves (2.18)--(2.19).

For the centered interval ([-r,r]), (7.1) becomes

\[
\max_A\sum_{i=1}^s\lambda_i(A)
=
%
r\max_{P\in\mathcal P_s(n)}|P|_{1,e}.
\tag{7.5}
\]

Thus even the optimal sum of the first \(s\) eigenvalues is exactly a finite-dimensional projection-constant problem.

\medskip\hrule\medskip

\subsection{8. Proof of the elementary variance bounds}

We use the following order-statistic lemma.

\paragraph{Lemma}

Let

\[
\mu_1\ge\cdots\ge\mu_n,
\qquad
\overline\mu=\frac1n\sum_i\mu_i,
\qquad
V=\sum_i(\mu_i-\overline\mu)^2.
\]

Then

\[
\mu_k
\le
\overline\mu+
\sqrt{\frac{n-k}{kn}V},
\tag{8.1}
\]

and

\[
\mu_k
\ge
\overline\mu-
\sqrt{\frac{k-1}{n(n-k+1)}V}.
\tag{8.2}
\]

\subparagraph{Proof}

Put \(\delta=\mu_k-\overline\mu\). If \(\delta\le0\), (8.1) is immediate. Otherwise the first \(k\) deviations are at least \(\delta\), and the sum of the remaining deviations is at most \(-k\delta\). Hence, by Cauchy--Schwarz,

\[
V
\ge
k\delta^2+\frac{k^2\delta^2}{n-k}
=
%
\frac{kn}{n-k}\delta^2.
\]

This proves (8.1). Formula (8.2) follows by applying (8.1) to
\(-\mu_n,\dots,-\mu_1\). \(\square\)

For a maximizer of \(U_k\), take all diagonal entries equal to \(b\). Then

\[
\operatorname{tr}A=nb
\]

and

\[
\begin{aligned}
\sum_{i=1}^n(\lambda_i(A)-b)^2
&=\operatorname{tr}(A-bI)^2\\
&=2\sum_{i<j}a_{ij}^2\\
&\le n(n-1)\rho^2.
\end{aligned}
\]

Applying (8.1) gives (2.14).

For a minimizer of \(L_k\), all diagonal entries equal \(a\), and the same argument with (8.2) gives (2.15).

\medskip\hrule\medskip

\subsection{9. Boundary cases and necessity of the hypotheses}

\subsubsection{9.1 The degenerate interval \(a=b=t\)}

There is only one matrix, \(tJ\). Its spectrum is

\[
{nt,0,\dots,0}.
\]

If \(t>0\), \(nt\) is the largest eigenvalue; if \(t<0\), it is the smallest. All formulas above reduce correctly to this spectrum.

\medskip\hrule\medskip

\subsubsection{9.2 Dimension \(n=2\)}

There are no genuinely intermediate eigenvalues. Formulas (2.1)--(2.4) give the complete answer for both eigenvalues.

For instance, if \(0\le a\le b\),

\[
{\lambda_2(A):A\in S_2([a,b])}
=[a-b,b-a].
\]

\medskip\hrule\medskip

\subsubsection{9.3 The sign assumption in the \(\lambda_2\) theorem is essential}

For \(n=5\), let \(G=C_5\), and put

\[
S=J-I-2A(G),\qquad M=I+S=J-2A(G).
\]

Then \(M\in S_5([-1,1])\), and its spectrum is

\[
1+\sqrt5,\quad 1+\sqrt5,\quad 1,\quad
1-\sqrt5,\quad 1-\sqrt5.
\]

Thus

\[
\lambda_2(M)=1+\sqrt5.
\]

The nonnegative-interval proof cannot be extended by simply substituting \(a<0\): negative cross-entries can reinforce, rather than suppress, the two sign parts of an eigenvector.

\medskip\hrule\medskip

\subsubsection{9.4 Endpoint matrices do not always suffice}

For the centered interval ([-1,1]),

\[
U_n(n;-1,1)=1,
\]

attained by \(I\). Its off-diagonal entries are interior points of the interval. Thus a brute-force enumeration of matrices with all entries in (\{-1,1\}) cannot solve Q706 in general.

This is one reason the density-matrix minimax formula (2.10) is needed.

\medskip\hrule\medskip

\subsection{10. Exact remaining obstruction}

The results above do \textbf{not} give closed formulas for every intermediate pair ((n,k)). Already in the centered case, the exact problem becomes

\[
\eta_{n,k}
=
%
\max_{\dim V=k}
\min_{\substack{X\succeq0,\ \operatorname{tr}X=1\\
\operatorname{Ran}X\subseteq V}}
|X|_{1,e}.
\tag{10.1}
\]

The Ky Fan relaxation replaces the inner minimum by the particular choice

\[
X=\frac1kP_V
\]

and gives

\[
\eta_{n,k}
\le
\frac1k
\max_{P\in\mathcal P_k(n)}|P|_{1,e},
\tag{10.2}
\]

but equality is not automatic: it would require the optimized endpoint matrix to have its first \(k\) eigenvalues equal on the relevant subspace.

The same projection geometry, together with Seidel matrices and equiangular tight frames, is central in current work on extremizing intermediate graph eigenvalues. Recent results still formulate this as a difficult family of extremal problems and obtain sharp bounds and equality classifications only in special ranks. ([arXiv][1])

A concrete next lemma would be:

\textbf{Target lemma.} Characterize the subspaces \(V\) for which the inner minimizer in (10.1) is a normalized orthogonal projection \(P_V/k\), and characterize when an entrywise endpoint matrix \(A\) satisfies
\[
P_VAP_V=\eta_{n,k}P_V.
\]
Such a result would convert the exact Ky Fan envelope into an exact formula for the individual \(k\)-th eigenvalue and reduce extremizers to structured two-eigenvalue sign or block matrices.

At present, (10.1) is the precise unresolved core left by this solution.

[1]: https://arxiv.org/html/2608.03196v1 "Extremal graphs for the 𝑘-th eigenvalue"
